\documentclass[12pt]{article}
\usepackage[utf8]{inputenc}
\usepackage[T1]{fontenc}
\usepackage{hyperref}
\title{Sunda Language Resources for K-12 Curriculum Development}
\author{Compiled by Claude Sonnet 5.5 (claude-sonnet-5-5) for Andr\'as Kornai}
\date{}
\begin{document}
\maketitle

\section{Introduction}
Sundanese (Basa Sunda, ISO 639-3 \texttt{sun}) is an Austronesian language of western Java, Indonesia. It is taught as a local-content (\emph{muatan lokal}) subject in schools. The sources found here are mostly Indonesian academic papers and bookseller pages, and they are not a full survey of what exists.

\section{Primary Educational Resources}
\begin{itemize}
\item \emph{Rancag\'e Diajar Basa Sunda}, a Grade 1 primary textbook approved as a local-content textbook in West Java and in circulation since 2016. A Telkom University thesis proposes a redesign. \url{https://openlibrary.telkomuniversity.ac.id/pustaka/files/153578/bab1/perancangan-ulang-buku-paket-pelajaran-bahasa-sunda-sekolah-dasar-kelas-1.pdf}
\item \emph{Klik! Diajar Basa Sunda} (Gramedia), a primary-level book aligned with Kurikulum 2013 (revised 2017). \url{https://www.gramedia.com/products/sdmi-klii-diajar-basa-sunda-jl2-k13-rev-2017}
\item Sundanese textbooks for senior high school grades X--XII, listed on Komunitas Belajar. \url{https://www.komunitasbelajar.id/2021/11/buku-bahasa-sunda-kelas-x-10-xi-11-xii.html}
\end{itemize}

\section{Linguistic and Reference Materials}
\begin{itemize}
\item \emph{Kamus Bahasa Sunda untuk Pelajar dan Umum} (Gramedia), a general dictionary for students and the public. It includes a section on speech levels (\emph{undak usuk basa}). \url{https://www.gramedia.com/products/kamus-bahasa-sunda-untuk-pelajar-dan-umum}
\end{itemize}

\section{Religious and Community Resources}
No religious or community resources were examined in the searches.

\section{Digital and Audio Resources}
No digital or audio resources were identified in the searches.

\section{Curriculum Development Context}
Sundanese is a local-content subject under Kurikulum 2013 and the newer Merdeka-era curriculum, with a West Java governor's regulation of 2014 encouraging regional-language instruction. Research papers report that student interest is low and that textbooks are widely seen as poorly designed and dull. The papers below discuss this.
\begin{itemize}
\item Telkom University journal article on Sundanese textbook design. \url{https://journals.telkomuniversity.ac.id/demandia/article/download/2160/1412}
\item Education journal article (UPI). \url{https://ejournal.upi.edu/index.php/edubasic/article/download/27483/pdf}
\end{itemize}
Current provincial rules and the list of approved textbooks should be confirmed with the West Java education office or the publishers.
\end{document}